\scanpage{36}
\mspage{35}

On a $A\otimes_{\mathbb Z_p}W$ étale sur $A$, ext. résiduelle
\[
k\otimes_{\mathbb F_p}k\simeq k^\Sigma,\qquad\Sigma=\operatorname{Aut}k=\operatorname{Aut}W,
\]
\[
A\otimes_{\mathbb Z_p}W\simeq A^\Sigma,\qquad \Sigma\simeq\mathbb Z/f\mathbb Z
\]
en correspondance bijection canonique avec $[0,f-1]$ par $\nu\mapsto\sigma^\nu$\quad\deleted{$(\unsure{\sigma\lambda=\lambda^p})$}.

On a
\[
\sigma_A(x_0,x_1,\ldots,x_{f-1})=(x_1,\ldots,x_{f-1},x_0).
\]
La donnée d'un module \inserted{\unsure{libre}} $M$ \inserted{sur $A\otimes_{\mathbb Z_p}W$} équivaut à la donnée de modules $(M_\nu)_{\nu\in\Sigma}$ sur $A$; le rang de $M$ est la somme des rangs des $M_\nu$ \inserted{(on le veut égal à $h$)}. La donnée d'un endom $F_M$ $\sigma_A$-linéaire équivaut à la donnée d'hom. $A$-linéaires
\diagram{\begin{tikzcd}
M_0 \arrow[r,"\varphi_0"] & M_1 \arrow[r,"\varphi_1"] & \cdots \arrow[r] & M_{f-1} \arrow[r,"\varphi_{f-1}"] & M_0
\end{tikzcd}}
La condition $(*)$ \deleted{\unsure{s'écrit}} implique que les $f$ composés circulaires sont des isogénies, donc les $\varphi_i$ sont injectifs, donc les $\varphi_i$ des isogénies (car les rangs des $M_\nu$ tous égaux $\ldots$). \deleted{Donc} --- Utilisant maintenant le fait que la somme des rangs des $M_\nu$ est $h$, on trouve que les $M_\nu$ sont inversibles.
\margin{ce serait aussi une \unsure{conséquence} de \deleted{$(*')$} $(*)$ \unsure{et de|avec} $\pi$ un élét non nul \uline{quelconque} de $A$}

\deleted{On peut alors trouver un unique module \unsure{inversible} $\Omega$ sur $A$,}
On peut alors identifier tous les $M_\nu$ à \deleted{$M_0=\unsure{\Omega}$,} $M_0=\unsure{\mathfrak t}$, \deleted{\unsure{grâce à}} en choisissant \deleted{\unsure{comme}} pour \unsure{tt $\nu$}, l'isom \inserted{$0\leq\nu\leq f-1$}
\[
M_0\overset\sim\longrightarrow M_{\unsure{i|\nu}}
\]
donné par
\[
\pi^{-m_\nu}(\varphi_{\nu-1}\varphi_{\nu-2}\cdots\varphi_0),
\]
où $m_\nu$ est un entier $\geq0$ uniquement déterminé $(0\leq m_0\leq m_1\leq\cdots\leq m_{f-1})$. On trouve donc un isom can.
\[
M\simeq\Omega\otimes_{\mathbb Z_p}W=\Omega\otimes_A(A\otimes_{\mathbb Z_p}W).
\]
